\documentclass[12pt]{article}
\usepackage[utf8]{inputenc}
\usepackage[T1]{fontenc}
\usepackage[margin=1in]{geometry}
\usepackage{setspace}
\usepackage{xcolor}
\usepackage[colorlinks=true, urlcolor=blue!60!black, linkcolor=blue!60!black]{hyperref}
\setlength{\parindent}{2em}
\setlength{\parskip}{1.25em}
\renewcommand{\baselinestretch}{1.0}
\begin{document}
\begin{center}
{\LARGE \textit{Planned Parenthood of Southeastern Pennsylvania v.\ Casey}, 505 U.S.\ 833 (1992)\par}
\vspace{0.8em}
{\large Case Brief\par}
\end{center}
\vspace{1em}

\section{Memory Jogger}

In a fractured decision authored jointly by Justices O'Connor, Kennedy, and Souter, the Supreme Court reaffirmed the ``essential holding'' of \textit{Roe v.\ Wade} on \textit{stare decisis} and Fourteenth Amendment liberty grounds, but replaced \textit{Roe}'s trimester framework and strict scrutiny with the ``undue burden'' standard---upholding four Pennsylvania abortion regulations while striking down spousal notification (later overruled by \textit{Dobbs v.\ Jackson Women's Health Organization}, 597 U.S.\ 215 (2022)).

\section{Detailed Case Facts}

In 1988 and 1989, the Pennsylvania legislature amended the Pennsylvania Abortion Control Act of 1982, 18 Pa.\ Cons.\ Stat.\ \S\S\ 3203--3220 (1990), adding or tightening five regulatory provisions governing abortion access in the Commonwealth (\href{https://www.courtlistener.com/opinion/112786/planned-parenthood-of-southeastern-pa-v-casey/}{CourtListener}; \href{https://supreme.justia.com/cases/federal/us/505/833/}{Justia}):
\begin{enumerate}
\item \textbf{Informed Consent and 24-Hour Waiting Period (\S\ 3205):} Required that, except in a medical emergency, a physician orally inform a woman at least 24 hours before an abortion of the nature of the procedure, the health risks of abortion and childbirth, and the ``probable gestational age of the unborn child.'' A physician or qualified nonphysician also had to inform her of the availability of state-published printed materials describing fetal development, medical assistance for childbirth, paternal child-support liability, and adoption alternatives, and the woman had to certify in writing that she had been informed of those materials.
\item \textbf{Parental Consent with Judicial Bypass (\S\ 3206):} Required an unemancipated minor under age 18 to obtain the informed consent of one parent (or guardian), unless a court of common pleas authorized the abortion after finding either that the minor was mature and capable of giving informed consent or that the abortion was in her best interests.
\item \textbf{Spousal Notification (\S\ 3209):} Prohibited a physician from performing an abortion on a married woman unless she provided a signed statement certifying that she had notified her husband, subject to four narrow exceptions: her husband was not the father; he could not be located after diligent effort; the pregnancy resulted from spousal sexual assault reported to law enforcement within 90 days; or she reasonably believed notification would cause her husband or someone else to inflict bodily injury upon her. Physicians who violated \S\ 3209 faced license revocation and civil liability to the husband, including \$5,000 in punitive damages.
\item \textbf{Medical Emergency Definition (\S\ 3203):} Excused compliance with the informed-consent, parental-consent, and spousal-notice requirements only when a physician determined in good faith clinical judgment that a condition necessitated an immediate abortion ``to avert her death or for which a delay will create serious risk of substantial and irreversible impairment of a major bodily function.''
\item \textbf{Recordkeeping and Reporting (\S\S\ 3207(b), 3214(a), 3214(f)):} Required abortion facilities to file facility-identification reports (made public if the facility received state-appropriated funds), quarterly aggregate reports by trimester, and confidential patient-anonymous individual reports for each abortion, including under \S\ 3214(a)(12) a married woman's reason for failing to notify her husband.
\end{enumerate}

Before any of the challenged amendments took effect, five abortion clinics (led by Planned Parenthood of Southeastern Pennsylvania) and one physician representing himself and a class of abortion-providing physicians sued Governor Robert P.\ Casey and other state officials, seeking declaratory and injunctive relief on the ground that all five provisions were facially unconstitutional under the Due Process Clause of the Fourteenth Amendment (\href{https://supreme.justia.com/cases/federal/us/505/833/}{Justia}).

\section{Procedural History}

\begin{itemize}
\item \textbf{U.S.\ District Court for the Eastern District of Pennsylvania (1990):} After issuing a preliminary injunction and conducting a three-day bench trial, the district court held all five challenged provisions facially unconstitutional under \textit{Roe v.\ Wade}'s strict-scrutiny and trimester framework, made extensive factual findings regarding the burdens of the 24-hour waiting period and the prevalence of domestic violence under \S\ 3209, and entered a permanent injunction against enforcement (\textit{Planned Parenthood of Southeastern Pennsylvania v.\ Casey}, 744 F.\ Supp.\ 1323 (E.D.\ Pa.\ 1990)) (\href{https://supreme.justia.com/cases/federal/us/505/833/}{Justia}).
\item \textbf{U.S.\ Court of Appeals for the Third Circuit (1991):} Concluding that intervening Supreme Court plurality opinions in \textit{Webster v.\ Reproductive Health Services}, 492 U.S.\ 490 (1989), and \textit{Hodgson v.\ Minnesota}, 497 U.S.\ 417 (1990), had displaced strict scrutiny in favor of Justice O'Connor's ``undue burden'' test, a three-judge panel affirmed in part and reversed in part: it upheld the medical emergency definition, informed consent and 24-hour waiting period, parental consent, and reporting requirements, but struck down the spousal notification requirement as an undue burden (\textit{Planned Parenthood of Southeastern Pennsylvania v.\ Casey}, 947 F.2d 682 (3d Cir.\ 1991)) (\href{https://law.justia.com/cases/federal/appellate-courts/F2/947/682/153995/}{Justia 3d Cir.}). Then-Circuit Judge Samuel A.\ Alito dissented in part, arguing that spousal notification should also be upheld.
\item \textbf{U.S.\ Supreme Court (1992):} Both sides petitioned for certiorari (Nos.\ 91-744 and 91-902), and the United States appeared as \textit{amicus curiae} urging the Court to overrule \textit{Roe v.\ Wade}. Oral argument took place on April 22, 1992, and the Court issued its decision on June 29, 1992, affirming the Third Circuit in No.\ 91-902 and affirming in part, reversing in part (as to \S\ 3214(a)(12)), and remanding in No.\ 91-744 (\href{https://www.law.cornell.edu/supct/html/91-744.ZS.html}{Cornell LII Syllabus}).
\item \textbf{Subsequent History:} Thirty years later, in \textit{Dobbs v.\ Jackson Women's Health Organization}, 597 U.S.\ 215 (2022) (authored by Justice Alito), the Supreme Court expressly overruled both \textit{Roe v.\ Wade} and \textit{Casey}.
\end{itemize}

\section{Judicial Votes}

The Court fractured sharply across five opinions, producing a complex part-by-part alignment (\href{https://www.law.cornell.edu/supct/html/91-744.ZS.html}{Cornell LII Syllabus}; \href{https://supreme.justia.com/cases/federal/us/505/833/}{Justia}):
\begin{itemize}
\item \textbf{Joint Opinion (Justices Sandra Day O'Connor, Anthony M.\ Kennedy, and David H.\ Souter):} Announced the judgment of the Court; delivered the \textbf{Opinion of the Court (5--4)} as to Parts I, II, III, V-A, V-C, and VI (joined by Justices Blackmun and Stevens); delivered a \textbf{four-Justice opinion} as to Part V-E (joined by Justice Stevens); and delivered a \textbf{three-Justice plurality opinion} as to Parts IV, V-B, and V-D.
\item \textbf{Justice John Paul Stevens} (concurring in part and dissenting in part): Joined Parts I, II, III, V-A, V-C, V-E, and VI; dissented from Parts IV, V-B, and V-D.
\item \textbf{Justice Harry A.\ Blackmun} (concurring in part, concurring in the judgment in part, and dissenting in part): Joined Parts I, II, III, V-A, V-C, and VI; concurred in the judgment invalidating \S\ 3214(a)(12); dissented from upholding \S\ 3205, \S\ 3206, and the remaining reporting provisions.
\item \textbf{Chief Justice William H.\ Rehnquist} (concurring in the judgment in part and dissenting in part), joined by \textbf{Justices Byron R.\ White, Antonin Scalia, and Clarence Thomas}: Concurred in the judgment upholding \S\ 3203, \S\ 3205, \S\ 3206, and the upheld reporting provisions; dissented from reaffirming \textit{Roe} and from invalidating \S\ 3209 and \S\ 3214(a)(12).
\item \textbf{Justice Antonin Scalia} (concurring in the judgment in part and dissenting in part), joined by \textbf{Chief Justice Rehnquist and Justices White and Thomas}: Same judgment alignment as Chief Justice Rehnquist, writing separately to reject any constitutional liberty interest in abortion and to attack the joint opinion's \textit{stare decisis} and undue-burden reasoning.
\end{itemize}

\noindent \textbf{Issue-by-Issue Vote Totals:}
\begin{itemize}
\item Reaffirming \textit{Roe}'s essential holding: \textbf{5--4} (O'Connor, Kennedy, Souter, Stevens, Blackmun vs.\ Rehnquist, White, Scalia, Thomas).
\item Medical emergency definition (\S\ 3203): \textbf{9--0} upheld (as construed by the Third Circuit).
\item Informed consent and 24-hour waiting period (\S\ 3205): \textbf{7--2} upheld (O'Connor, Kennedy, Souter, Rehnquist, White, Scalia, Thomas vs.\ Stevens, Blackmun).
\item One-parent consent with judicial bypass (\S\ 3206): \textbf{7--2} upheld (same 7--2 alignment).
\item Spousal notification (\S\ 3209): \textbf{5--4} invalidated (O'Connor, Kennedy, Souter, Stevens, Blackmun vs.\ Rehnquist, White, Scalia, Thomas).
\item Recordkeeping and reporting (\S\S\ 3207(b), 3214(a), 3214(f), excluding spousal-notice reporting): \textbf{8--1} upheld (all except Blackmun).
\item Spousal-notice reporting (\S\ 3214(a)(12)): \textbf{5--4} invalidated (same 5--4 alignment as \S\ 3209).
\end{itemize}

\section{Holding}

\textbf{Judgment:} The Supreme Court affirmed the Third Circuit's judgment in No.\ 91-902 (upholding the invalidation of \S\ 3209's spousal notification requirement) and affirmed in part, reversed in part, and remanded in No.\ 91-744 (upholding \S\ 3203, \S\ 3205, \S\ 3206, and the general reporting provisions of \S\S\ 3207(b), 3214(a), and 3214(f), while invalidating \S\ 3214(a)(12)'s spousal-notice reporting requirement) (\href{https://www.law.cornell.edu/supct/html/91-744.ZS.html}{Cornell LII Syllabus}).

\textbf{Rule Announced:} By a five-Justice majority (Parts I--III), the Court reaffirmed the three-part ``essential holding'' of \textit{Roe v.\ Wade}: (1) a woman has a Fourteenth Amendment liberty right to choose an abortion before fetal viability without undue state interference; (2) after viability, the State may restrict or proscribe abortion so long as it includes exceptions for pregnancies endangering the woman's life or health; and (3) from the outset of pregnancy, the State has legitimate interests in protecting maternal health and potential fetal life. By a controlling three-Justice plurality (Part IV, operating as the narrowest grounds supporting the judgment), the Court abandoned \textit{Roe}'s trimester framework and strict scrutiny in favor of the \textbf{undue burden standard}: a pre-viability regulation is unconstitutional if its ``purpose or effect is to place a substantial obstacle in the path of a woman seeking an abortion before the fetus attains viability.'' In facial challenges, a restriction is judged by its effect on the group for whom it is an actual restriction (the ``large fraction'' test), rather than all women seeking abortions (\href{https://supreme.justia.com/cases/federal/us/505/833/}{Justia}).

\textbf{Holding vs.\ Dicta:} The reaffirmation of \textit{Roe}'s viability line, the four-factor \textit{stare decisis} framework (workability, reliance, doctrinal erosion, and changed factual premises), the invalidation of \S\ 3209 and \S\ 3214(a)(12), the validation of \S\ 3203, \S\ 3205, \S\ 3206, and the remaining reporting provisions, and the partial overruling of \textit{Akron I} and \textit{Thornburgh} were necessary to the Court's judgment. By contrast, the joint opinion's broad philosophical statement in Part II that ``[a]t the heart of liberty is the right to define one's own concept of existence, of meaning, of the universe, and of the mystery of human life'' (later invoked in \textit{Lawrence v.\ Texas}, 539 U.S.\ 558 (2003)) and its extended commentary in Part III-C on the Court's institutional legitimacy when overruling ``watershed'' decisions ``under fire'' sweep far beyond the resolution of Pennsylvania's statute and function largely as influential dicta.

\section{Analysis of Opinions}

\subsection{Joint Opinion (Justices O'Connor, Kennedy, and Souter)}

\textbf{Substantive Due Process (Part II, Majority):} Grounding the abortion right in the ``liberty'' protected by the Fourteenth Amendment's Due Process Clause rather than a freestanding ``right of privacy,'' the majority rejected both strict incorporationism (limiting liberty to the Bill of Rights) and Justice Scalia's narrow historical-tradition test from \textit{Michael H.\ v.\ Gerald D.}, 491 U.S.\ 110 (1989). Quoting Justice John Marshall Harlan II's dissent in \textit{Poe v.\ Ullman}, 367 U.S.\ 497, 542--43 (1961), the Court described liberty as a ``rational continuum'' requiring ``reasoned judgment.'' Linking abortion to \textit{Griswold}, \textit{Eisenstadt}, \textit{Loving}, and bodily-integrity precedents, the majority stressed that the physical and personal burdens of pregnancy are ``unique to the human condition'' and cannot be conscripted by the State (\href{https://supreme.justia.com/cases/federal/us/505/833/}{Justia}).

\textbf{\textit{Stare Decisis} and Institutional Legitimacy (Part III, Majority):} Evaluating \textit{Roe} under four prudential factors, the majority concluded that: (1) \textit{Roe} had not proved unworkable; (2) two decades of societal reliance meant that ``[t]he ability of women to participate equally in the economic and social life of the Nation has been facilitated by their ability to control their reproductive lives''; (3) no doctrinal evolution had rendered \textit{Roe} an anachronism; and (4) medical advances moving viability earlier (from 28 weeks to 23--24 weeks) did not undermine viability itself as the constitutional dividing line. Distinguishing \textit{West Coast Hotel Co.\ v.\ Parrish} (overruling \textit{Lochner}/\textit{Adkins}) and \textit{Brown v.\ Board of Education} (overruling \textit{Plessy v.\ Ferguson}) as responses to fundamentally changed factual understandings, the majority warned that overruling a watershed decision ``under fire'' without special justification would inflict ``profound and unnecessary damage to the Court's legitimacy'' (\href{https://supreme.justia.com/cases/federal/us/505/833/}{Justia}).

\textbf{Undue Burden Framework and Application (Parts IV \& V):} In Part IV (plurality), the trio discarded \textit{Roe}'s trimester framework because it undervalued the State's profound interest in potential life before viability. States could henceforth enact persuasive informed-consent measures favoring childbirth so long as they did not impose a ``substantial obstacle.'' Applying this standard in Part V, the Court (with Blackmun and Stevens) deferred to the Third Circuit's broad reading of \S\ 3203's medical emergency exception (Part V-A) and struck down \S\ 3209's spousal notification requirement (Part V-C), citing the district court's domestic-violence findings and concluding that even if \S\ 3209 affected fewer than 1\% of all abortion patients, it operated as a substantial obstacle in a ``large fraction'' of the cases where it was relevant (married women who did not wish to notify their husbands). Relying on \textit{Danforth} and rejecting common-law coverture, the majority declared that ``[w]omen do not lose their constitutionally protected liberty when they marry.'' In Parts V-B and V-D (plurality), the trio overruled \textit{Akron I} and \textit{Thornburgh} in part to uphold truthful, nonmisleading informed consent, the 24-hour waiting period (finding the record insufficient for facial invalidation despite added travel and cost), and one-parent consent with judicial bypass. In Part V-E (joined by Stevens), the Court upheld general reporting requirements while striking down \S\ 3214(a)(12)'s spousal-notice reporting mandate (\href{https://supreme.justia.com/cases/federal/us/505/833/}{Justia}).

\subsection{Concurring in Part and Dissenting in Part (Justice Stevens)}

Justice Stevens joined Parts I, II, III, V-A, V-C, V-E, and VI, praising the majority's \textit{stare decisis} analysis and emphasizing that under \textit{Roe}, a fetus is not a Fourteenth Amendment ``person'' and the State's interest in potential life must remain secular. He dissented from Parts IV, V-B, and V-D, arguing that while states may require neutral medical-risk disclosures (\S\S\ 3205(a)(1)(i) and (iii)), they may not inject persuasive, state-sponsored anti-abortion messaging (\S\S\ 3205(a)(2)(i)--(iii)) or force adult women to endure a mandatory 24-hour waiting period. Even under the undue burden test, Stevens contended, a waiting period rested on ``outmoded and unacceptable assumptions about the decisionmaking capacity of women'' and lacked any legitimate justification (\href{https://supreme.justia.com/cases/federal/us/505/833/}{Justia}).

\subsection{Concurring in Part, Concurring in the Judgment in Part, and Dissenting in Part (Justice Blackmun)}

Justice Blackmun---author of \textit{Roe v.\ Wade}---hailed the joint opinion as ``an act of personal courage and constitutional principle'' that preserved \textit{Roe}'s flame, and joined Parts I, II, III, V-A, V-C, and VI. However, he argued that strict scrutiny and the trimester framework remained the proper constitutional standard, invoking both bodily integrity and Fourteenth Amendment gender equality (noting that abortion restrictions ``conscript[] women's bodies into [the State's] service''). Under strict scrutiny, Blackmun would have invalidated the informed-consent counseling, 24-hour waiting period, parental-consent, and reporting provisions alike, concurring only in the judgment striking down \S\ 3214(a)(12). He closed with a sobering personal warning that the constitutional right to abortion survived by a single vote (\href{https://www.law.cornell.edu/supct/html/91-744.ZX2.html}{Cornell LII Blackmun Opinion}).

\subsection{Concurring in the Judgment in Part and Dissenting in Part (Chief Justice Rehnquist, joined by Justices White, Scalia, and Thomas)}

Chief Justice Rehnquist argued that \textit{Roe} was wrongly decided and should be overruled, contending that abortion is \textit{sui generis} because it involves the purposeful termination of potential life and therefore cannot be equated with \textit{Griswold}, \textit{Loving}, or \textit{Pierce}. He criticized the joint opinion's ``newly minted variation on \textit{stare decisis},'' noting the irony of professing fidelity to \textit{Roe} while discarding its trimester framework and overruling \textit{Akron I} and \textit{Thornburgh}. Calling the undue burden test unguided and ``created largely out of whole cloth,'' Rehnquist advocated rational-basis review under which every provision of Pennsylvania's statute---including spousal notification, which rationally furthered a husband's interest in procreation and marital integrity---should be upheld (\href{https://supreme.justia.com/cases/federal/us/505/833/}{Justia}).

\subsection{Concurring in the Judgment in Part and Dissenting in Part (Justice Scalia, joined by Chief Justice Rehnquist and Justices White and Thomas)}

Justice Scalia went further than Rehnquist by asserting that abortion is not a constitutionally protected ``liberty'' at all because ``(1) the Constitution says absolutely nothing about it, and (2) the longstanding traditions of American society have permitted it to be legally proscribed.'' He ridiculed the joint opinion's ``reasoned judgment'' as a ``collection of adjectives'' masking personal value judgments, dismissed the undue burden standard as a ``verbal shell game,'' and mocked the majority's ``keep what you want and throw away the rest'' version of \textit{stare decisis}. Most famously, Scalia condemned the Court's institutional-legitimacy rationale as a ``Nietzschean vision'' of an ``Imperial Judiciary,'' comparing \textit{Roe} and \textit{Casey} to \textit{Dred Scott v.\ Sandford} and arguing that judicial entrenchment of abortion rights poisoned democratic politics and judicial confirmation hearings (\href{https://www.law.cornell.edu/supct/html/91-744.ZX4.html}{Cornell LII Scalia Opinion}).

\section{Examples: Future Applications}

\begin{enumerate}
\item \textbf{Same-Side (Unconstitutional Undue Burden):} A state enacts a law requiring every physician who performs abortions to maintain active admitting privileges at a local hospital within 30 miles, where hospitals routinely deny privileges to abortion providers for non-medical reasons, forcing half of the state's clinics to close without demonstrable health benefits. \textit{Why it falls here:} Under \textit{Casey}'s undue burden standard, unnecessary health regulations that have the purpose or effect of placing a substantial obstacle in the path of women seeking pre-viability abortions---drastically increasing travel distances and wait times without conferring meaningful medical benefits---are unconstitutional (cf.\ \textit{Whole Woman's Health v.\ Hellerstedt}, 579 U.S.\ 582 (2016)).
\item \textbf{Same-Side (Unconstitutional Undue Burden):} A state requires a married woman seeking a first-trimester medication abortion to obtain her husband's written acknowledgment unless she produces a police report proving felony domestic battery. \textit{Why it falls here:} Like Pennsylvania's \S\ 3209, this law directly empowers an abusive spouse to exert an effective veto over a woman's pre-viability decision, and its narrow police-report exception fails to protect women who fear retaliation, psychological coercion, or harm to their children.
\item \textbf{Opposite-Side (Constitutional Under \textit{Casey}):} A state requires abortion clinics to provide patients, 24 hours before a non-emergency abortion, with an accurate, non-misleading state brochure detailing gestational development, medical risks of abortion and childbirth, and free prenatal and adoption services, while permitting the initial consultation to occur via telehealth so only one in-person clinic visit is needed. \textit{Why it falls here:} Because \textit{Casey} expressly upheld Pennsylvania's truthful informed-consent disclosures and 24-hour waiting period as permissible expressions of the State's interest in potential life, an even less burdensome telehealth version clearly imposes no substantial obstacle.
\item \textbf{Opposite-Side (Constitutional Under \textit{Casey}):} A state bans elective abortions after 26 weeks of gestation (past fetal viability) while expressly exempting abortions that a physician, in good faith medical judgment, deems necessary to preserve the pregnant woman's life or physical or mental health. \textit{Why it falls here:} \textit{Casey}'s second core holding explicitly confirms that after viability, the State's interest in potential human life becomes sufficiently compelling to justify proscribing abortion, provided life and health exceptions are preserved.
\item \textbf{Fence-Sitter (Close / Fact-Dependent Case):} A state imposes a 72-hour mandatory waiting period requiring two separate in-person trips to a clinic in a largely rural state with only two operating abortion facilities, resulting in average delays of two weeks and substantial lost wages and lodging costs for low-income patients. \textit{Why it is a fence-sitter:} \textit{Casey} upheld Pennsylvania's 24-hour waiting period on a facial challenge despite recognizing ``troubling'' travel and cost burdens, stating that a ``particular burden is not of necessity a substantial obstacle'' on that record; yet \textit{Casey} left open whether a longer delay combined with a stronger empirical record of actual clinic closures and thwarted access would cross the line into an unconstitutional substantial obstacle.
\end{enumerate}

\section{Critique}

\textbf{Originalist and Textualist Critiques:} Originalist scholars and jurists sharply attacked \textit{Casey} on three fronts. First, echoing Justice Scalia's dissent, critics argue that the Fourteenth Amendment's Due Process Clause contains no textual or historical warrant for an unenumerated right to abortion, given widespread nineteenth-century state prohibitions at ratification (\href{https://www.law.cornell.edu/supct/html/91-744.ZX4.html}{Cornell LII Scalia Opinion}). Second, scholars such as Michael Stokes Paulsen have condemned \textit{Casey}'s ``super-\textit{stare decisis}'' doctrine as an illegitimate claim of judicial supremacy that entrenches constitutional error precisely because the error is controversial (\href{https://www.thepublicdiscourse.com/2012/06/5772/}{Paulsen, Public Discourse}). Third, conservative critics contend that the ``undue burden'' test lacked objective legal standards, inviting ad hoc judicial policymaking by district judges while leaving the viability line philosophically arbitrary.

\textbf{Progressive, Feminist, and Doctrinal Critiques:} Progressive and feminist scholars also voiced deep reservations. First, then-Judge Ruth Bader Ginsburg and scholars such as Reva Siegel argued that anchoring reproductive freedom solely in substantive due process ``privacy'' or ``liberty''---rather than the Equal Protection Clause and women's equal citizenship---left abortion rights doctrinally vulnerable, even while welcoming \textit{Casey}'s acknowledgment of women's reliance interests and anti-coverture principles (\href{https://www.uclalawreview.org/the-equality-argument-for-abortion-rights/}{UCLA Law Review}; \href{https://texaslawreview.org/wp-content/uploads/2023/04/Siegel.pdf}{Siegel, Texas Law Review}). Second, reproductive-rights scholars criticized the plurality's abandonment of strict scrutiny and its facial approval of Pennsylvania's 24-hour waiting period and biased-counseling rules, warning that the malleable ``undue burden'' standard invited a wave of targeted state restrictions that disproportionately harmed poor, rural, and minor women (\href{https://scholarship.law.wm.edu/cgi/viewcontent.cgi?article=1055&context=wmjowl}{William \& Mary Journal of Women and the Law}).

\textbf{Competing Underlying Values:} At its core, \textit{Casey} reflects an intractable collision among four constitutional values: (1) individual bodily autonomy and women's equal participation in economic and social life; (2) the State's asserted moral interest in protecting prenatal life; (3) democratic self-governance and federalism on deeply contested moral questions; and (4) rule-of-law stability, reliance, and the institutional legitimacy of an unelected judiciary.

\section{Key Quotations}

\begin{quote}
``Liberty finds no refuge in a jurisprudence of doubt.''
\end{quote}
\noindent --- Joint Opinion of Justices O'Connor, Kennedy, and Souter (Opinion of the Court, Part I), 505 U.S.\ at 844 (\href{https://supreme.justia.com/cases/federal/us/505/833/}{Justia}; \href{https://www.courtlistener.com/opinion/112786/planned-parenthood-of-southeastern-pa-v-casey/}{CourtListener}).

\begin{quote}
``At the heart of liberty is the right to define one's own concept of existence, of meaning, of the universe, and of the mystery of human life. Beliefs about these matters could not define the attributes of personhood were they formed under compulsion of the State.''
\end{quote}
\noindent --- Joint Opinion of Justices O'Connor, Kennedy, and Souter (Opinion of the Court, Part II), 505 U.S.\ at 851 (\href{https://supreme.justia.com/cases/federal/us/505/833/}{Justia}).

\begin{quote}
``An undue burden exists, and therefore a provision of law is invalid, if its purpose or effect is to place a substantial obstacle in the path of a woman seeking an abortion before the fetus attains viability.''
\end{quote}
\noindent --- Joint Opinion of Justices O'Connor, Kennedy, and Souter (Plurality Opinion, Part IV), 505 U.S.\ at 878 (\href{https://supreme.justia.com/cases/federal/us/505/833/}{Justia}).

\begin{quote}
``I am 83 years old. I cannot remain on this Court forever, and when I do step down, the confirmation process for my successor well may focus on the issue before us today. That, I regret, may be exactly where the choice between the two worlds will be made.''
\end{quote}
\noindent --- Justice Blackmun (concurring in part, concurring in the judgment in part, and dissenting in part), 505 U.S.\ at 943 (\href{https://www.law.cornell.edu/supct/html/91-744.ZX2.html}{Cornell LII}).

\begin{quote}
``The States may, if they wish, permit abortion on demand, but the Constitution does not require them to do so. The permissibility of abortion, and the limitations upon it, are to be resolved like most important questions in our democracy: by citizens trying to persuade one another and then voting.''
\end{quote}
\noindent --- Justice Scalia (concurring in the judgment in part and dissenting in part), 505 U.S.\ at 979 (\href{https://www.law.cornell.edu/supct/html/91-744.ZX4.html}{Cornell LII}).

\end{document}
